% Generated by publication/build_sections.py; do not edit by hand.
\paragraph{Informal verdict.}
\textbf{A --- Complete solution}, with \textbf{high confidence} in the two requested regimes.

\paragraph{Lean coverage.}
The exact 17-file retry-02 closure unconditionally proves the fixed-alpha all-n sure/superpolynomial/limit regime and, for every fixed n >= 2, the exact kappa\_n upper-edge relative error and both ratio limits. The arbitrary-n product law, attainment, measurability, small-ball, localization, slicing, and ellipsoid inputs are discharged internally. The complete aggregate passed independent pinned Lean CI.

\subsection{Audited informal answer}
Q788 --- rigorous solution
\begin{enumerate}
\item Faithful restatement
\end{enumerate}
Identify the unit circle \(\mathbb U\subset\mathbb R^2\) with

\[
\mathbb T:=\{z\in\mathbb C:|z|=1\}.
\]

Let \(Z_1,\ldots,Z_n\) be independent Haar-uniform random points of \(\mathbb T\), and put

\[
P_n(z):=\prod_{j=1}^n(z-Z_j),\qquad
D_n:=\max_{|z|=1}|P_n(z)|.
\]

Thus \(D_n\) is the maximum, over \(M\in\mathbb U\), of the product of the Euclidean chord distances \(MM_j\).
The problem states that \(D_n\in[2,2^n]\) almost surely and asks for estimates of

\[
\mathbb P(D_n\ge \alpha)
\]

in the two regimes:


\(\alpha\) fixed and \(n\to\infty\);


\(n\) fixed and \(\alpha\uparrow 2^n\).


There is one evident typographical omission in the PDF: in the sentence
\(\mathbb P(D\in[2,2^n])=1\), the variable \(D\) should be \(D_n\). There is no other substantive ambiguity. Luc Pommeret

\begin{enumerate}
\item Probability space and main result
\end{enumerate}
Work on

\[
\Omega=[0,2\pi)^{\mathbb N},\qquad
\mathcal F=\mathcal B([0,2\pi))^{\otimes\mathbb N},\qquad
\mathbb P=\left(\frac{d\theta}{2\pi}\right)^{\otimes\mathbb N}.
\]

Let \(\Theta_j\) be the coordinate maps and \(Z_j=e^{i\Theta_j}\). Then the \(Z_j\) are independent and Haar-uniform on \(\mathbb T\).
The map

\[
(\theta_1,\ldots,\theta_n)\longmapsto
\max_{t\in[0,2\pi]}
\prod_{j=1}^n|e^{it}-e^{i\theta_j}|
\]

is continuous, since it is the maximum over a compact set of a jointly continuous function. Hence \(D_n\) is a well-defined random variable.
Define

\[
\kappa_n:=
\frac{\sqrt n}{\Gamma\!\left(\frac{n+1}{2}\right)}
\left(\frac2\pi\right)^{\frac{n-1}{2}}.
\]

Theorem
Let \(n\ge2\).
\(a\) Fixed \(\alpha\), \(n\to\infty\)
If \(\alpha\le2\), then

\[
\mathbb P(D_n\ge\alpha)=1.
\]

If \(\alpha>2\), then for every \(A>0\) there is a constant
\(C_{\alpha,A}<\infty\) such that

\[
\boxed{\;
\mathbb P(D_n<\alpha)\le C_{\alpha,A}n^{-A}.
\;}
\]

Consequently,

\[
\boxed{\;
\mathbb P(D_n\ge\alpha)
 =1-O_{\alpha,A}(n^{-A})
 \qquad\text{for every }A>0.
\;}
\]

There is also a complementary explicit lower estimate: for every fixed
\(\alpha>2\), there is \(c_\alpha>0\) such that, for all sufficiently large \(n\),

\[
\boxed{\;
\mathbb P(D_n<\alpha)
 \ge n!\left(\frac{c_\alpha}{\pi n\log n}\right)^n
 \ge \exp\!\bigl(-n\log\log n-O_\alpha(n)\bigr).
\;}
\]

Thus the exceptional probability tends to zero faster than every negative power of \(n\).
The natural, nondegenerate scale is exponential in \(\sqrt n\): there is a strictly positive random variable \(I^*\), expressible as the maximum of a Brownian-bridge functional, such that

\[
\boxed{\;
\frac{2\log D_n}{\sqrt n}\ \xrightarrow{\ d\ }\ I^*.
\;}
\]

Hence, at every continuity point \(x\ge0\) of the law of \(I^*\),

\[
\mathbb P\!\left(D_n\ge e^{x\sqrt n/2}\right)
 \longrightarrow \mathbb P(I^*\ge x).
\]

\(b\) Fixed \(n\), \(\alpha\uparrow2^n\)
Put

\[
\delta:=1-\frac{\alpha}{2^n}
       =\frac{2^n-\alpha}{2^n}.
\]

Then, as \(\delta\downarrow0\),

\[
\boxed{\;
\mathbb P\!\left(D_n\ge2^n(1-\delta)\right)
=
\kappa_n\,\delta^{\frac{n-1}{2}}
 \bigl(1+O_n(\delta)\bigr).
\;}
\]

Equivalently,

\[
\boxed{\;
\mathbb P(D_n\ge\alpha)
\sim
\frac{\sqrt n}{\Gamma\!\left(\frac{n+1}{2}\right)}
\left(\frac2\pi\right)^{\frac{n-1}{2}}
\left(\frac{2^n-\alpha}{2^n}\right)^{\frac{n-1}{2}}.
\;}
\]

Moreover,

\[
\mathbb P(D_n=2^n)=0.
\]


\begin{enumerate}
\item Fixed-level estimates
\end{enumerate}
3.1 A deterministic Fourier consequence of \(D_n\le\alpha\)
For a deterministic configuration \(Z_j=e^{i\theta_j}\), define

\[
u(t):=\log|P_n(e^{it})|
     =\sum_{j=1}^n\log|e^{it}-e^{i\theta_j}|.
\]

The logarithmic singularities are integrable.
We use the classical \(L^1\)-Fourier expansion

\[
\log|1-e^{ix}|
 =\log\!\left(2\left\lvert{}\sin\frac x2\right\rvert{}\right)
 =-\sum_{r=1}^{\infty}\frac{\cos(rx)}r.
\]

For completeness, this follows by first writing, for \(0<\rho<1\),

\[
\log|1-\rho e^{ix}|
 =-\sum_{r=1}^{\infty}\frac{\rho^r\cos(rx)}r,
\]

and then letting \(\rho\uparrow1\) in \(L^1([0,2\pi])\). The singularity is bounded by an integrable multiple of \(1+|\log|x||\).
In particular,

\[
\frac1{2\pi}\int_0^{2\pi}
 \log|e^{it}-e^{i\theta}|\,dt=0.
\]

Therefore

\[
\frac1{2\pi}\int_0^{2\pi}u(t)\,dt=0.
\]

Suppose now that \(D_n\le\alpha\). Then \(u(t)\le\log\alpha\). Writing
\(u=u_+-u_-\), the zero-mean identity gives

\[
\int u_+=\int u_-,
\]

and hence

\[
\int_0^{2\pi}|u(t)|\,dt
 =2\int_0^{2\pi}u_+(t)\,dt
 \le4\pi\log\alpha.
\]

Let

\[
S_r:=\sum_{j=1}^n Z_j^r.
\]

The \(r\)-th Fourier coefficient of \(u\), for \(r\ge1\), is

\[
\widehat u(r)
 =-\frac1{2r}\sum_{j=1}^ne^{-ir\theta_j}
 =-\frac{\overline{S_r}}{2r}.
\]

Consequently,

\[
\frac{|S_r|}{2r}
 =|\widehat u(r)|
 \le\frac1{2\pi}\int_0^{2\pi}|u(t)|\,dt
 \le2\log\alpha.
\]

We have proved:
Lemma 1
For every \(r\ge1\),

\[
\boxed{\quad
D_n\le\alpha
\quad\Longrightarrow\quad
|S_r|\le4r\log\alpha.
\quad}
\]

Thus a uniformly bounded maximum forces every fixed collection of power sums to lie in a fixed bounded region.

3.2 A fixed-dimensional small-ball estimate
Lemma 2
Fix an integer \(K\ge1\) and positive constants \(R_1,\ldots,R_K\). Then there is \(C=C(K,R_1,\ldots,R_K)\) such that

\[
\mathbb P\bigl(|S_r|\le R_r\text{ for }1\le r\le K\bigr)
 \le Cn^{-K}.
\]

Proof
Introduce the \(2K\)-dimensional random vector

\[
Y(\theta):=
(\cos\theta,\sin\theta,\ldots,\cos K\theta,\sin K\theta)
 \in\mathbb R^{2K}.
\]

Then

\[
(\Re S_1,\Im S_1,\ldots,\Re S_K,\Im S_K)
 =\sum_{j=1}^nY(\Theta_j).
\]

Let

\[
\varphi(\xi)
 :=\mathbb E\bigl[e^{i\xi\cdot Y(\Theta_1)}\bigr]
 =\frac1{2\pi}\int_0^{2\pi}
 e^{i\xi\cdot Y(\theta)}\,d\theta
\]

be the characteristic function.
Behavior near the origin
Orthogonality of the trigonometric system gives

\[
\mathbb EY=0,\qquad
\mathbb E(YY^{\mathsf T})=\frac12I_{2K}.
\]

Taylor expansion therefore yields

\[
\varphi(\xi)
 =1-\frac14|\xi|^2+O_K(|\xi|^3).
\]

Hence there are \(c,\rho>0\), depending only on \(K\), such that

\[
|\varphi(\xi)|\le e^{-c|\xi|^2},
\qquad |\xi|\le\rho.
\]

Behavior away from the origin
If \(|\varphi(\xi)|=1\), equality in the triangle inequality for the defining integral implies that

\[
e^{i\xi\cdot Y(\theta)}
\]

is constant for almost every \(\theta\), hence everywhere. Thus the trigonometric polynomial
\(\xi\cdot Y(\theta)\) is constant modulo \(2\pi\), and by continuity is constant as a real function. Since it has no constant Fourier coefficient, it must vanish identically. Linear independence of the sine and cosine functions then gives \(\xi=0\). Therefore

\[
|\varphi(\xi)|<1\qquad(\xi\ne0).
\]

We also need integrability at infinity. Write \(\xi=Rv\), where
\(R=|\xi|\) and \(|v|=1\), and let

\[
p_v(\theta):=v\cdot Y(\theta).
\]

This is a nonzero real trigonometric polynomial of degree at most \(K\).
There is a constant \(c_K>0\) such that

\[
\max_{1\le q\le2K+1}|p_v^{(q)}(\theta)|\ge c_K
\]

for every unit vector \(v\) and every \(\theta\). Indeed, otherwise compactness would give a nonzero \(p_v\) whose first \(2K+1\) derivatives vanish at some point. Then
\(p_v-p_v(\theta_0)\), after multiplication by \(e^{iK\theta}\), would give a polynomial of degree at most \(2K\) with a zero of order greater than \(2K\), forcing \(p_v\) to be constant, hence zero, a contradiction.
Partitioning the circle into a number of intervals depending only on \(K\), on each of which one of these derivatives stays bounded away from zero, the one-dimensional van der Corput estimate gives

\[
|\varphi(\xi)|
 \le C_K(1+|\xi|)^{-1/(2K+1)}.
\]

Indeed, on an interval where
\(|p_v^{(q)}|\ge c_K/2\), van der Corput gives

\[
\left\lvert{}\int e^{iRp_v(\theta)}\,d\theta\right\rvert{}
 \le C_{K}R^{-1/q}
 \le C_KR^{-1/(2K+1)}.
\]

Choose an integer

\[
N_0>2K(2K+1).
\]

Then \(|\varphi|^{N_0}\in L^1(\mathbb R^{2K})\).
Since \(|\varphi|<1\) away from zero and tends to zero at infinity, there is \(q\in(0,1)\) such that

\[
\sup_{|\xi|\ge\rho}|\varphi(\xi)|\le q.
\]

For \(n\ge N_0\),

\[
\begin{aligned}
\int_{\mathbb R^{2K}}|\varphi(\xi)|^n\,d\xi
&\le
 \int_{|\xi|<\rho}e^{-cn|\xi|^2}\,d\xi
 +q^{n-N_0}
   \int_{|\xi|\ge\rho}|\varphi(\xi)|^{N_0}\,d\xi\\
&\le C_Kn^{-K}+C_Kq^n
 \le C'_Kn^{-K}.
\end{aligned}
\]

Thus the \(n\)-fold convolution of the law of \(Y\) has a bounded continuous density

\[
f_n(x)=\frac1{(2\pi)^{2K}}
 \int_{\mathbb R^{2K}}e^{-i\xi\cdot x}\varphi(\xi)^n\,d\xi,
\]

and

\[
\|f_n\|_\infty\le C_Kn^{-K}.
\]

Finally, the set

\[
B=\{(w_1,\ldots,w_K)\in\mathbb C^K:
      |w_r|\le R_r\}
\]

has finite \(2K\)-dimensional volume. Hence

\[
\mathbb P\bigl((S_1,\ldots,S_K)\in B\bigr)
 \le \operatorname{vol}(B)\|f_n\|_\infty
 \le Cn^{-K}.
\]

This proves the lemma. \(\square\)

3.3 Superpolynomial convergence in part \(a\)
Fix \(\alpha>2\). By Lemma 1,

\[
\{D_n<\alpha\}
 \subseteq
 \bigcap_{r=1}^K\{|S_r|\le4r\log\alpha\}.
\]

Lemma 2 therefore gives

\[
\mathbb P(D_n<\alpha)\le C_{\alpha,K}n^{-K}.
\]

Given \(A>0\), choose an integer \(K>A\). Then

\[
\mathbb P(D_n<\alpha)=O_{\alpha,A}(n^{-A}),
\]

which proves the asserted upper estimate for the exceptional probability.
No limit and expectation have been interchanged here: the only limiting operation used in the probabilistic estimate was Fourier inversion after proving that \(\varphi^n\in L^1\).

3.4 A constructive lower bound for the exceptional probability
We now show that \(D_n<\alpha\) is not impossible and give an explicit lower estimate.
Let

\[
\omega_j=e^{2\pi ij/n},\qquad 0\le j<n,
\]

be the regular \(n\)-gon.
Lemma 3: stability of the regular polygon
There is an absolute constant \(C>0\) with the following property. Suppose

\[
Z_j=e^{i(2\pi j/n+\varepsilon_j)},\qquad
|\varepsilon_j|\le\eta.
\]

Then

\[
D_n\le
(2+n\eta)\exp\!\bigl(Cn\eta\log(en)\bigr).
\]

Proof
Fix \(z=e^{it}\) and put

\[
a_j:=|z-\omega_j|.
\]

Choose \(j_0\) so that \(a_{j_0}\) is minimal. Since

\[
|Z_j-\omega_j|\le|\varepsilon_j|\le\eta,
\]

we have

\[
|z-Z_j|\le a_j+\eta.
\]

For \(j\ne j_0\),

\[
a_j+\eta\le a_j e^{\eta/a_j}.
\]

Consequently,

\[
\prod_{j=0}^{n-1}|z-Z_j|
 \le
(a_{j_0}+\eta)
\prod_{j\ne j_0}a_j
\exp\!\left(\eta\sum_{j\ne j_0}\frac1{a_j}\right).
\]

Now

\[
\prod_{j=0}^{n-1}a_j=|z^n-1|\le2,
\]

while

\[
\prod_{j\ne j_0}a_j
=
\left\lvert{}\frac{z^n-1}{z-\omega_{j_0}}\right\rvert{}
=
\left\lvert{}
 \sum_{\ell=0}^{n-1}z^{n-1-\ell}\omega_{j_0}^{\,\ell}
\right\rvert{}
\le n,
\]

with the quotient interpreted by continuity when \(z=\omega_{j_0}\). Hence

\[
(a_{j_0}+\eta)\prod_{j\ne j_0}a_j\le2+n\eta.
\]

It remains to estimate the inverse distances. Since \(j_0\) is a nearest grid point, the grid points at cyclic distance \(k\) from \(j_0\) have angular distance at least
\((2k-1)\pi/n\) from \(t\). Using

\[
2\sin(x/2)\ge\frac{2x}{\pi},
\qquad 0\le x\le\pi,
\]

we obtain

\[
\sum_{j\ne j_0}\frac1{a_j}
 \le Cn\sum_{k=1}^{n}\frac1k
 \le Cn\log(en).
\]

Combining these estimates and maximizing over \(z\in\mathbb T\) proves the lemma. \(\square\)
Fix \(\alpha>2\). Choose \(c_\alpha>0\) sufficiently small that

\[
2e^{2Cc_\alpha}<\alpha.
\]

Set

\[
\eta_n:=\frac{c_\alpha}{n\log n}.
\]

Lemma 3 gives, for all sufficiently large \(n\),

\[
D_n\le
\left(2+\frac{c_\alpha}{\log n}\right)
\exp\!\left(
 Cc_\alpha\frac{\log(en)}{\log n}
\right)
<\alpha.
\]

Consider the \(n\) disjoint arcs of angular half-width \(\eta_n\) centered at the regular \(n\)-gon points. The probability that the \(n\) labelled random points occupy these arcs with exactly one point in each is

\[
n!\left(\frac{\eta_n}{\pi}\right)^n
=
n!\left(\frac{c_\alpha}{\pi n\log n}\right)^n.
\]

On this event the roots may be relabelled so that Lemma 3 applies. Thus

\[
\mathbb P(D_n<\alpha)
 \ge n!\left(\frac{c_\alpha}{\pi n\log n}\right)^n.
\]

Using \(n!\ge(n/e)^n\),

\[
\mathbb P(D_n<\alpha)
 \ge
\left(\frac{c_\alpha}{e\pi\log n}\right)^n
=
\exp\!\bigl(-n\log\log n-O_\alpha(n)\bigr).
\]


3.5 The actual probabilistic scale
A known theorem for this precise random-polynomial model gives a sharper description on the scale on which \(D_n\) fluctuates. For

\[
P_N(z)=\prod_{k=1}^N(z-z_k)^{m_k},
\qquad
s_N^2=\sum_{k=1}^Nm_k^2,
\]

the theorem assumes

\[
\sum_{k=1}^N
 \exp\!\left(-\varepsilon\frac{s_N}{m_k}\right)
 \longrightarrow0
\qquad(\varepsilon>0)
\]

and concludes that

\[
\frac1{s_N}\log\max_{|z|=1}|P_N(z)|^2
 \xrightarrow{\ d\ }I^*,
\]

where \(I^*>0\) almost surely and is obtained from a Brownian bridge.
Here \(m_k=1\), so \(s_N=\sqrt N\), and the condition becomes

\[
N e^{-\varepsilon\sqrt N}\longrightarrow0.
\]

It is therefore satisfied, giving

\[
\frac{2\log D_n}{\sqrt n}\xrightarrow{\ d\ }I^*.
\]

This is the theorem of Tucci and Whiting, and its hypotheses match the present model exactly. arXiv+1
This external theorem is not needed for the preceding superpolynomial fixed-\(\alpha\) estimate.

\begin{enumerate}
\item The upper-edge asymptotic
\end{enumerate}
We now fix \(n\ge2\) and let

\[
\alpha=2^n(1-\delta),\qquad \delta\downarrow0.
\]

4.1 Reduction by rotation
Because \(D_n\) is unchanged when all roots are multiplied by the same unit complex number, define

\[
x_1:=0,\qquad
x_j:=\operatorname{Arg}(Z_j\overline{Z_1})\in(-\pi,\pi],
\quad 2\le j\le n.
\]

The torus transformation

\[
(\Theta_1,\ldots,\Theta_n)
\longmapsto
(\Theta_1,\Theta_2-\Theta_1,\ldots,\Theta_n-\Theta_1)
\]

preserves Haar measure. Hence

\[
x_2,\ldots,x_n
\]

are independent and uniform on \((-\pi,\pi]\), independently of the discarded global rotation.
Set \(x=(x_2,\ldots,x_n)\in\mathbb R^{n-1}\).

4.2 Localization near a coincident configuration
Suppose

\[
D_n\ge2^n(1-\delta).
\]

Then for some \(z\in\mathbb T\),

\[
\prod_{j=1}^n\frac{|z-Z_j|}{2}\ge1-\delta.
\]

Every factor belongs to \([0,1]\). Therefore each factor is itself at least \(1-\delta\).
Let \(-e^{ic}\) denote \(z\), so that \(e^{ic}\) is its antipode. If \(d_j\in[0,\pi]\) is the circular angular distance between \(Z_j\) and \(e^{ic}\), then

\[
\frac{|z-Z_j|}{2}=\cos\frac{d_j}{2}.
\]

Thus

\[
\cos\frac{d_j}{2}\ge1-\delta.
\]

Using, for \(0\le y\le\pi/2\),

\[
1-\cos y\ge\frac{2y^2}{\pi^2},
\]

we obtain

\[
d_j=O(\sqrt\delta).
\]

All roots therefore lie in an arc of length \(O(\sqrt\delta)\). In particular,

\[
|x|=O_n(\sqrt\delta).
\]

This also shows that, for sufficiently small \(\delta\), there is no ambiguity from the choice of angular lifts.

4.3 Quadratic approximation
For \(|u|<\pi\), define

\[
h(u):=-\log\cos\frac u2.
\]

Near zero,

\[
h(u)=\frac{u^2}{8}+O(u^4).
\]

For \(x_1=0\) and \(x=(x_2,\ldots,x_n)\) sufficiently small,

\[
\frac{D_n}{2^n}
=
\max_{c\in\mathbb R}
 \prod_{j=1}^n
 \cos\frac{x_j-c}{2}.
\]

Hence

\[
\Phi(x):=-\log\frac{D_n}{2^n}
=
\min_{c\in\mathbb R}
 \sum_{j=1}^n h(x_j-c).
\]

Let

\[
\bar x:=\frac1n\sum_{j=1}^n x_j
       =\frac1n\sum_{j=2}^n x_j
\]

and define

\[
Q(x):=\min_{c\in\mathbb R}\sum_{j=1}^n(x_j-c)^2
     =\sum_{j=1}^n(x_j-\bar x)^2.
\]

The minimizing \(c\) for the function involving \(h\) lies between the smallest and largest \(x_j\), because \(h'\) is odd and strictly increasing near zero. It is therefore \(O(|x|)\). The Taylor expansion of \(h\), evaluated at this minimizer and at \(c=\bar x\), gives the two-sided estimate

\[
\boxed{\;
\Phi(x)=\frac18Q(x)+O_n(|x|^4).
\;}
\]

Since \(Q\) is positive definite on the variables \(x_2,\ldots,x_n\), and its smallest eigenvalue is \(1/n\), we have

\[
|x|^2\le nQ(x).
\]

Consequently the remainder may equivalently be written

\[
\Phi(x)=\frac18Q(x)+O_n(Q(x)^2).
\]


4.4 The relevant ellipsoid
Writing \(m=n-1\), one has

\[
Q(x)
 =\sum_{j=2}^n x_j^2
  -\frac1n\left(\sum_{j=2}^nx_j\right)^2
 =x^{\mathsf T}Ax,
\]

where

\[
A=I_m-\frac1nJ_m
\]

and \(J_m\) is the all-ones matrix.
The eigenvalues of \(A\) are


\(1\), with multiplicity \(m-1\);


\(1/n\), in the direction \((1,\ldots,1)\).


Therefore

\[
\det A=\frac1n.
\]

The event in question is

\[
\Phi(x)\le-\log(1-\delta)
         =\delta+O(\delta^2).
\]

The localization already proved gives \(Q(x)=O_n(\delta)\). Substituting this into the quadratic approximation shows that there is \(C_n<\infty\) such that, for all sufficiently small \(\delta\),

\[
\{Q(x)\le8\delta-C_n\delta^2\}
 \subseteq
\left\{\frac{D_n}{2^n}\ge1-\delta\right\}
 \subseteq
\{Q(x)\le8\delta+C_n\delta^2\}.
\]

Thus the desired probability is the normalized volume of an ellipsoid, up to relative error \(O_n(\delta)\).
Let

\[
v_m=\frac{\pi^{m/2}}{\Gamma(m/2+1)}
\]

be the volume of the Euclidean unit ball in \(\mathbb R^m\). Then

\[
\operatorname{vol}\{x:x^{\mathsf T}Ax\le8\delta\}
=
\frac{v_m(8\delta)^{m/2}}{\sqrt{\det A}}
=
\sqrt n\,v_m(8\delta)^{m/2}.
\]

Because the relative-angle vector is uniform with density \((2\pi)^{-m}\),

\[
\begin{aligned}
\mathbb P\!\left(D_n\ge2^n(1-\delta)\right)
&=
\frac{\sqrt n\,v_m8^{m/2}}{(2\pi)^m}
 \delta^{m/2}\bigl(1+O_n(\delta)\bigr)\\
&=
\frac{\sqrt n}{\Gamma(m/2+1)}
 \left(\frac2\pi\right)^{m/2}
 \delta^{m/2}\bigl(1+O_n(\delta)\bigr).
\end{aligned}
\]

Since \(m=n-1\), this is exactly

\[
\mathbb P\!\left(D_n\ge2^n(1-\delta)\right)
=
\frac{\sqrt n}{\Gamma\!\left(\frac{n+1}{2}\right)}
\left(\frac2\pi\right)^{\frac{n-1}{2}}
\delta^{\frac{n-1}{2}}
\bigl(1+O_n(\delta)\bigr).
\]

This proves part \(b\).

4.5 The exact endpoint
Every factor \(|z-Z_j|\) is at most \(2\). Therefore \(D_n=2^n\) can hold only if there is some \(z\in\mathbb T\) such that

\[
|z-Z_j|=2
\]

for every \(j\). This is equivalent to

\[
Z_1=\cdots=Z_n=-z.
\]

Conversely, coincident roots do give \(D_n=2^n\).
Since the \(Z_j\) are independent and have nonatomic distributions,

\[
\mathbb P(Z_1=\cdots=Z_n)=0,
\]

so

\[
\mathbb P(D_n=2^n)=0.
\]


\begin{enumerate}
\item Boundary and small-\(n\) checks
\end{enumerate}
5.1 Exact calculation for \(n=2\)
Let \(\Delta\in[0,\pi]\) be the smaller angular separation of the two roots. Since the angular difference is uniform, \(\Delta\) is uniform on \([0,\pi]\).
After rotating the roots to \(e^{\pm i\Delta/2}\),

\[
\begin{aligned}
D_2
&=\max_t
 |e^{it}-e^{i\Delta/2}|
 |e^{it}-e^{-i\Delta/2}|\\
&=2\max_t|\cos(\Delta/2)-\cos t|\\
&=2\bigl(1+\cos(\Delta/2)\bigr)\\
&=4\cos^2(\Delta/4).
\end{aligned}
\]

Hence, for \(2\le\alpha\le4\),

\[
\boxed{\;
\mathbb P(D_2\ge\alpha)
 =\frac4\pi\arccos\!\left(\frac{\sqrt\alpha}{2}\right).
\;}
\]

As \(\alpha\uparrow4\),

\[
\mathbb P(D_2\ge\alpha)
 \sim\frac2\pi\sqrt{4-\alpha}.
\]

Our general formula gives

\[
\kappa_2
=
\frac{\sqrt2}{\Gamma(3/2)}
 \left(\frac2\pi\right)^{1/2}
 =\frac4\pi,
\]

and, since \(\delta=(4-\alpha)/4\),

\[
\kappa_2\delta^{1/2}
 =\frac4\pi\frac{\sqrt{4-\alpha}}2
 =\frac2\pi\sqrt{4-\alpha},
\]

in exact agreement.
5.2 The first higher-dimensional case
For \(n=3\),

\[
\kappa_3
 =\frac{\sqrt3}{\Gamma(2)}\frac2\pi
 =\frac{2\sqrt3}{\pi}.
\]

Thus

\[
\mathbb P(D_3\ge8(1-\delta))
=
\frac{2\sqrt3}{\pi}\delta+O(\delta^2).
\]

The linear power is consistent with the fact that, after quotienting out global rotation, two independent relative angles must fall into a two-dimensional ellipsoid of radius \(O(\sqrt\delta)\).

\begin{enumerate}
\item Why the hypotheses matter
\end{enumerate}
The independence assumption is essential. Let \(U\) be uniform and set

\[
Z_j=e^{i(U+2\pi j/n)},\qquad 0\le j<n.
\]

Each \(Z_j\) is individually uniform, but the points are dependent. Their polynomial is

\[
P_n(z)=z^n-e^{inU},
\]

so

\[
D_n=2
\]

identically. Thus, for every fixed \(\alpha>2\),

\[
\mathbb P(D_n\ge\alpha)=0,
\]

the opposite of the independent result.
Uniformity is also essential. If the \(Z_j\) are independent but all have the degenerate distribution \(Z_j=1\), then

\[
D_n=2^n
\]

identically. In particular, both the \(\sqrt n\)-scale law and the upper-edge asymptotic fail.
Finally, \(n\ge2\) is relevant at the endpoint. For \(n=1\),

\[
D_1=2
\]

deterministically, whereas for every \(n\ge2\) the event \(D_n=2^n\) has probability zero.

Final answer
For every fixed \(\alpha>2\),

\[
\boxed{\;
\mathbb P(D_n\ge\alpha)
 =1-O_{\alpha,A}(n^{-A})
 \quad\text{for every }A>0,
\;}
\]

with the explicit complementary construction

\[
\mathbb P(D_n<\alpha)
 \ge
\exp\!\bigl(-n\log\log n-O_\alpha(n)\bigr).
\]

The typical size is \(D_n=\exp(\Theta_{\mathrm{law}}(\sqrt n))\).
For fixed \(n\) and \(\alpha\uparrow2^n\),

\[
\boxed{\;
\mathbb P(D_n\ge\alpha)
\sim
\frac{\sqrt n}{\Gamma\!\left(\frac{n+1}{2}\right)}
\left(\frac2\pi\right)^{\frac{n-1}{2}}
\left(\frac{2^n-\alpha}{2^n}\right)^{\frac{n-1}{2}}.
\;}
\]
